\subsection*{§. 198.\\ Clifford 平行性。}

椭圆几何中滑移的奇特性质，使 Clifford 引入了这种几何中两条直线平行的重要概念。Clifford 关于这个问题的思想，尤其由 Klein 介绍并发展\srcfn{1}{参见《Math. Annalen》第37卷中的论文及 Klein 的石印讲义：\emph{Nicht-Euklidische Geometrie}，Göttingen，1890。欲知更多详情者可参阅这些著作。}。若像欧几里得几何和双曲几何那样，把有一个公共无穷远点的两条直线定义为平行线，就必须说椭圆空间没有（实）无穷远点，因而没有平行线。然而，双曲几何的平行线只部分具有欧几里得几何平行线的性质；许多最简单的性质实际上都失效了。例如在欧几里得几何中，给定两条平行线，可以移动整个空间，使这两条平行线分别沿自身滑动；两条平行线有无穷多条公共垂线，等等。这些性质对双曲空间中的两条平行线不再成立。反之，在椭圆空间中采用 Clifford 的下述平行定义，就能重新得到这些性质：

\emph{椭圆空间中的两条直线，若都与绝对形的同一对共轭母线相交，即属于同一个 Clifford 直线汇，就称为平行。}

立即可见，与欧几里得几何和双曲几何不同，Clifford 意义下的两条平行线不在同一平面内。

平行性必须分为两类：右旋平行与左旋平行；从空间中任一点向一条给定直线可作两条平行线，一条为右旋，另一条为左旋。只有当所取点到给定直线的距离为一个象限 $\frac\pi2R$ 时，它们才重合（重合于给定直线的极线）。显然，给定 Clifford 意义下的两条平行线，存在一个空间滑移，使这两条直线分别沿自身滑动。若从一点 $A'$ 作线段，该点位于第二条直线 $a'$ 上，并令此线段 $A'A$ 在第一条直线的点 $A$ 处垂直于第一条直线，那么在连续滑移中，线段 $A'A$ 的两个端点将分别沿直线 $a,a'$ 滑动；由于它长度不变，并在所有位置都垂直于 $a$，也必垂直于 $a'$\srcfn{1}{在以直线 $AA'$ 为母线的直纹面上，$A$ 的轨迹确实是母线的正交轨线；因而由 Gauss 定理，直线 $a'$，即 $A'$ 的轨迹，也有同样性质。}。

因此，正如在欧几里得几何中一样，一条直线上的所有点到其每一条平行线的距离相等，而两条直线之间有无穷多个彼此相等的最短距离。

再注意，若对任意直线 $r$ 作一个滑移，例如右旋滑移；该滑移不使它沿自身滑动，而把它移到新位置 $r'$，则直线 $r,r'$ 是左旋平行的。事实上，在右旋滑移中，绝对形上第二族的所有母线分别沿自身滑动；两条（共轭）母线 $g_1,g_2$ 也如此，而 $r$ 正与它们相交；所以 $r'$ 也与 $g_1,g_2$ 相交，因而在左旋意义下与 $r$ 平行。由此可知，若两条平行线 $r,r'$ 被同一直线 $p$ 截交，则它们与 $p$ 在两个交点 $A,B$ 处形成的（内错）角相等，正如欧几里得几何中一样。事实上，若 $r,r'$ 的平行性为右旋，就作一个把 $A$ 移到 $B$ 的左旋滑移；$p$ 将沿自身滑动，而 $r$ 与 $r'$ 重合。
